
les
flèches écrites ici sont des homomorphismes de restriction, et on utilise le fait que \(i^*\) est un isomorphisme d'après le théorème de changement de base propre (SGA 4 XII 5.5). Je dis que \emph{le carré} (7.13.2) \emph{est commutatif}. La démonstration est essentiellement la même que celle qui établit la commutativité de (7.10.2); il suffit ici d'utiliser la compatibilité de la formation des classes de Chern avec les images inverses, dans le cas des morphismes \(j:X_t\to X\) et \(i:X_s\to X\).

\textbf{7.13.3.} Notons que le carré (7.13.2) provient par passage à la limite projective des carrés analogues, où la cohomologie à coefficients dans le faisceau \(\ell\)-adique \(\mathbb Z_\ell(i)\) est remplacée par celle à coefficients dans les faisceaux de torsion \(\mu_n^{\otimes i}\) (où \(n\) parcourt les puissances de \(\ell\)). Donc la commutativité de (7.13.2) ne fait qu'exprimer la commutativité des carrés faisant intervenir les classes de Chern mod. \(n\), pour toute puissance \(n\) de \(\ell\). Choisissons une clôture algébrique \(\bar K\) de \(K=k(t)\), d'où une clôture algébrique \(\bar k\) de \(k=k(s)\) (savoir le corps résiduel de la clôture normale de \(V\) dans \(\bar K\)). On peut, pour toute sous-extension finie \(K_\alpha\) de \(\bar K\), considérer le diagramme commutatif analogue (pour \(n\) fixé) relatif à la situation déduite de \(f:X\to S\) par le changement de base \(S_\alpha\to S\), \(S_\alpha\) étant le normalisé de \(S\) dans \(K_\alpha\). On notera pour ceci que, par Krull--Akizuki, et \(S\) étant hensélien, \(S_\alpha\) est bien un trait; de plus, il faudra supposer ici que les fibres de \(X\) sur \(S\) sont géométriquement régulières, i.e. \(X\) \emph{lisse} sur \(S\), de sorte que cette hypothèse se conserve par changement de base. On trouve ainsi, pour \(\alpha\) variable, un système inductif de carrés commutatifs. Passant à la limite inductive et utilisant (SGA 4 VII 5.8) et (IV 3.2), on trouve un carré commutatif (cf. 7.16) :
